\section{Robustness}
\label{sec:robustness}

\subsection{Marginal Impact by LL Experience}
\label{sec:robustness_marginal}

The main analysis restricts to first-time LL candidates. A natural question is whether subsequent LL opportunities have similar effects, or whether the marginal impact diminishes with experience. To test this, I return to the full (unrestricted) estimation sample and stratify by the number of prior LL awards: $n = 0$ (first LL, matching the main results), $n = 1$ (second LL), and $n \geq 2$ (third or more). Within each experience group, I compare treated players ($D_{ie} = 1$) to controls ($D_{ie} = 0$) who have the same number of prior LL awards, using the full pre-treatment control specification.

Tables~\ref{tab:marginal_gs} and~\ref{tab:marginal_ngs} report stratum-by-stratum estimates. These are descriptive: each stratum is estimated on its own subsample, the Grand Slam $n = 1$ and $n \geq 2$ cells are small (ATP $N = 48$ and $N = 80$; WTA $N = 31$ and $N = 19$), and the resulting point estimates are unstable, in places exceeding the physically possible treatment dose in absolute value. Section~\ref{sec:robustness_saturated} estimates the experience gradient on the pooled sample in a single specification, which is the basis for the paper's claims about how the return varies with $n$.

\begin{table}[htbp]
\centering
\caption{Marginal Impact by LL Experience: Grand Slam Samples}
\label{tab:marginal_gs}
\begin{threeparttable}
\scriptsize
\begin{tabular}{l l*{5}{c}}
\toprule
Tour & Experience & 4w & 8w & 12w & 26w & 52w \\
\midrule
ATP & First LL ($n=0$) & 72.5$^{***}$ & 72.1$^{***}$ & 82.8$^{***}$ & 72.6$^{**}$ & 23.5 \\
  &  & (26.7) & (26.6) & (26.1) & (36.2) & (59.1) \\
  &  & -4.4 & 0.2 & -4.8 & 4.7 & -9.7 \\
  &  & (9.1) & (10.9) & (10.3) & (13.5) & (21.0) \\
  &  & -0.2 & 0.1 & 0.1 & 0.7 & 0.8 \\
  &  & (0.5) & (0.5) & (0.5) & (0.5) & (1.5) \\
  &  & 0.3 & 0.7 & 0.9 & 1.5 & 1.3 \\
  &  & (1.0) & (1.0) & (0.9) & (1.1) & (2.8) \\
[0.3em]
ATP & Second LL ($n=1$) & -220.9$^{***}$ & -215.3$^{***}$ & -195.6$^{***}$ & -153.7$^{**}$ & -202.1$^{**}$ \\
  &  & (48.5) & (44.9) & (49.0) & (64.1) & (98.6) \\
  &  & -97.9$^{***}$ & -89.9$^{***}$ & -96.5$^{***}$ & -83.7$^{***}$ & -98.6$^{***}$ \\
  &  & (17.6) & (18.3) & (21.3) & (23.0) & (26.7) \\
  &  & -1.6$^{**}$ & -1.7$^{***}$ & -1.2$^{**}$ & 0.6 & 1.3 \\
  &  & (0.6) & (0.5) & (0.5) & (0.6) & (1.9) \\
  &  & -9.5$^{***}$ & -9.3$^{***}$ & -8.4$^{***}$ & -4.5$^{***}$ & -2.9 \\
  &  & (1.4) & (1.1) & (1.2) & (1.5) & (3.9) \\
[0.3em]
ATP & Third+ LL ($n \geq 2$) & -30.7 & -38.3 & -40.4 & 4.3 & -14.8 \\
  &  & (40.6) & (35.8) & (35.7) & (41.6) & (62.6) \\
  &  & -8.3 & -15.0 & -18.6 & -13.4 & -26.5 \\
  &  & (15.2) & (13.6) & (14.3) & (20.5) & (20.1) \\
  &  & 0.4 & 0.7 & 0.8 & 0.5 & 0.7 \\
  &  & (0.8) & (0.7) & (0.7) & (0.8) & (1.7) \\
  &  & 0.9 & 1.0 & 0.9 & 0.7 & 1.8 \\
  &  & (1.4) & (1.2) & (1.2) & (1.3) & (2.6) \\
[0.3em]
\midrule
WTA & First LL ($n=0$) & 40.0 & 55.0$^{**}$ & 27.8 & 7.5 & 26.5 \\
  &  & (28.5) & (22.5) & (20.1) & (42.8) & (86.4) \\
  &  & 13.3$^{*}$ & 15.0$^{*}$ & 11.8 & 21.6$^{*}$ & 36.8$^{*}$ \\
  &  & (7.6) & (8.5) & (9.1) & (12.7) & (18.9) \\
  &  & -0.1 & 0.3 & 0.3 & 0.7 & 1.8 \\
  &  & (0.5) & (0.5) & (0.6) & (0.7) & (1.5) \\
  &  & -0.5 & 0.8 & 0.6 & 1.9$^{*}$ & 3.2 \\
  &  & (1.0) & (0.7) & (0.8) & (1.0) & (2.5) \\
[0.3em]
WTA & Second LL ($n=1$) & 54.6 & 33.0 & -82.8$^{**}$ & -154.7$^{***}$ & -181.3$^{*}$ \\
  &  & (61.9) & (53.2) & (38.9) & (49.4) & (104.2) \\
  &  & -6.1 & -16.8 & -17.9 & -16.3 & -52.4 \\
  &  & (10.0) & (10.0) & (14.8) & (18.8) & (41.5) \\
  &  & -0.6 & -0.2 & -0.5 & -0.7 & -2.4 \\
  &  & (1.2) & (1.2) & (1.1) & (1.3) & (2.2) \\
  &  & -1.0 & -0.4 & -0.9 & -1.7 & -6.0 \\
  &  & (2.1) & (2.0) & (1.9) & (2.1) & (3.9) \\
[0.3em]
\bottomrule
\end{tabular}

\begin{tablenotes}
\footnotesize
\item \textit{Notes.} Full (unrestricted) Grand Slam sample stratified by number of prior LL awards ($n$). Within each group, treated = received LL, control = eligible but not selected, both with the same number of prior awards. Event FE specification with full pre-treatment controls (including prior-LL count variables). Standard errors clustered at the player level. $^{***}p<0.01$, $^{**}p<0.05$, $^{*}p<0.10$.
\end{tablenotes}
\end{threeparttable}
\end{table}

\begin{table}[htbp]
\centering
\caption{Marginal Impact by LL Experience: Non-Grand-Slam Samples}
\label{tab:marginal_ngs}
\begin{threeparttable}
\scriptsize
\begin{tabular}{l l*{5}{c}}
\toprule
Tour & Experience & 4w & 8w & 12w & 26w & 52w \\
\midrule
ATP & First LL ($n=0$) & 20.1$^{**}$ & 19.7$^{*}$ & 13.5 & 44.5$^{***}$ & 37.7 \\
  &  & (10.2) & (10.4) & (10.3) & (14.5) & (30.4) \\
  &  & 6.8 & 11.1$^{*}$ & 10.0$^{*}$ & 28.0$^{**}$ & 30.4 \\
  &  & (6.3) & (5.7) & (6.0) & (12.7) & (20.5) \\
  &  & -1.4$^{***}$ & -0.9$^{***}$ & -0.3$^{*}$ & 1.7$^{***}$ & 5.2$^{***}$ \\
  &  & (0.2) & (0.2) & (0.2) & (0.3) & (0.7) \\
  &  & -2.6$^{***}$ & -1.7$^{***}$ & -0.7$^{**}$ & 3.0$^{***}$ & 9.0$^{***}$ \\
  &  & (0.4) & (0.4) & (0.3) & (0.5) & (1.2) \\
[0.3em]
ATP & Second LL ($n=1$) & 6.1 & 11.4 & 2.6 & 19.6 & -1.1 \\
  &  & (24.1) & (26.1) & (28.1) & (34.7) & (53.4) \\
  &  & -3.5 & -1.5 & -1.5 & -6.2 & -10.9 \\
  &  & (7.9) & (8.6) & (9.1) & (10.4) & (14.3) \\
  &  & -1.3$^{***}$ & -0.8$^{**}$ & -0.2 & 1.5$^{***}$ & 4.2$^{***}$ \\
  &  & (0.4) & (0.4) & (0.4) & (0.5) & (1.0) \\
  &  & -2.4$^{***}$ & -1.4$^{*}$ & -0.6 & 2.7$^{***}$ & 7.4$^{***}$ \\
  &  & (0.8) & (0.8) & (0.8) & (1.0) & (2.0) \\
[0.3em]
ATP & Third+ LL ($n \geq 2$) & 11.6 & 11.6 & 15.5 & 19.1 & -25.5 \\
  &  & (22.8) & (22.7) & (22.9) & (25.5) & (43.2) \\
  &  & 4.2 & 2.8 & 3.7 & -2.0 & -2.3 \\
  &  & (7.1) & (7.4) & (7.8) & (9.4) & (13.3) \\
  &  & -0.9$^{**}$ & -0.4 & -0.1 & 1.3$^{***}$ & 3.2$^{***}$ \\
  &  & (0.4) & (0.3) & (0.3) & (0.4) & (1.0) \\
  &  & -1.4$^{*}$ & -0.6 & 0.0 & 2.4$^{***}$ & 5.9$^{***}$ \\
  &  & (0.7) & (0.7) & (0.7) & (0.9) & (1.9) \\
[0.3em]
\midrule
WTA & First LL ($n=0$) & 17.8 & 18.8$^{*}$ & 23.0$^{*}$ & 20.5 & 59.2 \\
  &  & (11.3) & (11.0) & (12.2) & (17.9) & (39.6) \\
  &  & 6.5$^{*}$ & 7.3$^{*}$ & 10.8$^{***}$ & 6.4 & 5.7 \\
  &  & (3.6) & (3.8) & (4.1) & (5.8) & (8.1) \\
  &  & -1.1$^{***}$ & -0.7$^{***}$ & -0.3$^{**}$ & 0.9$^{***}$ & 2.7$^{***}$ \\
  &  & (0.2) & (0.1) & (0.1) & (0.2) & (0.5) \\
  &  & -2.1$^{***}$ & -1.4$^{***}$ & -0.8$^{***}$ & 1.6$^{***}$ & 5.6$^{***}$ \\
  &  & (0.3) & (0.3) & (0.3) & (0.4) & (1.1) \\
[0.3em]
WTA & Second LL ($n=1$) & 36.6 & 29.6 & 34.4 & 20.9 & -18.0 \\
  &  & (27.9) & (29.3) & (29.3) & (36.4) & (65.0) \\
  &  & 8.0 & 5.2 & 6.2 & 4.9 & 9.4 \\
  &  & (6.5) & (6.9) & (7.4) & (9.5) & (11.8) \\
  &  & -1.0$^{***}$ & -0.6 & -0.1 & 1.4$^{***}$ & 3.4$^{***}$ \\
  &  & (0.4) & (0.4) & (0.4) & (0.5) & (0.9) \\
  &  & -1.6$^{**}$ & -1.0 & -0.0 & 2.6$^{***}$ & 6.2$^{***}$ \\
  &  & (0.7) & (0.7) & (0.7) & (0.9) & (1.8) \\
[0.3em]
WTA & Third+ LL ($n \geq 2$) & 17.2 & -1.5 & -17.1 & -81.4$^{**}$ & -159.6$^{**}$ \\
  &  & (26.3) & (29.0) & (29.8) & (38.9) & (72.1) \\
  &  & 3.3 & 1.4 & 3.6 & -1.4 & -24.9$^{*}$ \\
  &  & (6.4) & (7.1) & (6.9) & (10.0) & (15.0) \\
  &  & -0.4 & -0.0 & 0.4 & 1.1$^{***}$ & 2.4$^{**}$ \\
  &  & (0.3) & (0.3) & (0.3) & (0.4) & (1.1) \\
  &  & -1.0 & -0.6 & -0.4 & 0.5 & 2.7 \\
  &  & (0.7) & (0.7) & (0.7) & (1.0) & (2.1) \\
[0.3em]
\bottomrule
\end{tabular}

\begin{tablenotes}
\footnotesize
\item \textit{Notes.} Full (unrestricted) non-Grand-Slam sample stratified by number of prior LL awards ($n$). Control function specification with $\hat{v}_{ie} \times \text{horizon}$ terms and full pre-treatment controls. Standard errors clustered at the player level. $^{***}p<0.01$, $^{**}p<0.05$, $^{*}p<0.10$.
\end{tablenotes}
\end{threeparttable}
\end{table}

\subsection{The Experience Gradient}
\label{sec:robustness_saturated}

Comparing coefficients estimated on separate stratum subsamples does not identify how the return varies with experience, because each subsample carries its own covariate distribution and its own residual variance. A pooled regression coefficient is a variance-weighted average of stratum effects whose implicit weights depend on the within-stratum variance of treatment conditional on covariates, and those weights need not match the population share of each stratum \citep{Sloczynski2022_interpreting_ols}. Differences between a pooled coefficient and a share-weighted average of subsample coefficients are therefore uninformative about heterogeneity.

I instead estimate the gradient in a single saturated specification on the pooled sample:
\begin{equation}
\label{eq:saturated}
\begin{split}
Y_{ie,h} = {} & \alpha_e + \alpha_h + \sum_h \beta_h D_{ie}\mathbf{1}[h] \\
 & + \sum_h \gamma_{1h} D_{ie}\mathbf{1}[h]\mathbf{1}[n_{ie}=1]
   + \sum_h \gamma_{2h} D_{ie}\mathbf{1}[h]\mathbf{1}[n_{ie}\geq 2] \\
 & + \sum_h \lambda_{1h} \mathbf{1}[h]\mathbf{1}[n_{ie}=1]
   + \sum_h \lambda_{2h} \mathbf{1}[h]\mathbf{1}[n_{ie}\geq 2] \\
 & + \Gamma Z^{pre}_{ie} + \varepsilon_{ie,h},
\end{split}
\end{equation}
where $n_{ie}$ counts prior LL awards. Here $\beta_h$ is the effect for first-time recipients, $\beta_h + \gamma_{1h}$ the effect for second-time recipients, and $\beta_h + \gamma_{2h}$ the effect for third-plus recipients. The $\lambda_{jh}$ terms let each experience stratum carry its own untreated outcome path across horizons; without them the specification would force LL-naive and repeat candidates to share a common untreated trajectory, and any difference would be absorbed into $\hat\gamma_{jh}$ and misread as treatment-effect heterogeneity. All coefficients come from one regression on one sample, so they aggregate to the pooled estimand by construction and the question becomes a hypothesis test rather than an accounting identity.

The specification must be saturated on the control side as well as the treatment side. Without stratum-specific horizon effects, the model forces the untreated outcome path to share a common horizon profile across experience strata, and since players who reach $n \geq 2$ have by construction spent years losing in final qualifying, their untreated ranking trajectory differs from a first-timer's. Any such difference would be absorbed into $\hat\gamma_{2h}$ and misread as a treatment shifter. Equation~(\ref{eq:saturated}) therefore includes $\mathbf{1}[n=1]$ and $\mathbf{1}[n \geq 2]$ interacted with the horizon indicators, so each stratum carries its own untreated path and $\hat\gamma_{jh}$ is identified from the treated-versus-control contrast within stratum.

Table~\ref{tab:saturated_experience} reports the results, and they support a weaker claim than the paper's original framing supposed. On the ATP Grand Slam sample the first-time effect is 55.2 ranking points at 4 weeks, 58.4 at 8 weeks, and 69.1 at 12 weeks, each significant at the 1\% level. The second-LL shifter is small and statistically indistinguishable from zero at every horizon ($-6.7$, $3.7$, and $55.4$ at 4, 12, and 26 weeks, none approaching significance), so a player receiving a second lucky loser slot gains about as much as a player receiving a first. The third-plus shifter is negative and significant at the three short horizons ($-79.7$, $p = 0.014$; $-87.4$, $p = 0.011$; $-99.3$, $p = 0.007$), enough to offset the base effect, but it is indistinguishable from zero at 26 and 52 weeks.

Crucially, the joint Wald test does \textit{not} reject a flat experience profile for ranking points on any of the four samples (ATP Grand Slam $F = 1.20$, $p = 0.285$; non-Grand-Slam ATP $F = 0.54$, $p = 0.861$; WTA Grand Slam $F = 1.21$, $p = 0.281$; non-Grand-Slam WTA $F = 1.08$, $p = 0.374$). Only main draws entered on the ATP Grand Slam sample rejects ($F = 2.15$, $p = 0.019$). The honest summary is therefore that the return to an opportunity shows no statistically established gradient in the number already received, with the qualified exception of short-horizon ranking points and main-draw access on the ATP Grand Slam sample, where third-plus recipients do worse.

What the evidence does rule out is the specific claim that the \textit{first} opportunity is distinctive. The second-LL shifter is small and insignificant on every sample and at every horizon, and the implied second-LL effects on the ATP Grand Slam sample (48.5, 52.8, 72.8 at 4, 8, and 12 weeks) track the first-LL effects closely. Whatever gradient exists does not begin between the first opportunity and the second. The restriction to first-time recipients in the main analysis remains the right sample choice for estimating the effect of an access shock cleanly, but it should not be read as isolating a uniquely pivotal event.

Two caveats bound this evidence. First, $n_{ie}$ is itself an outcome of prior treatment: within a stratum the lottery randomizes $D$, so $\beta_h$ and $\beta_h + \gamma_{jh}$ are each valid effects for their own population, but the comparison across strata is between endogenously constituted groups, and the joint test is a test of equality across treatment-history-defined populations rather than of a causal dose-response. Second, an earlier version of this specification omitted the stratum-by-horizon terms and produced a $\hat\gamma_{2h}$ that grew monotonically with the horizon, reaching $-182.3$ at 52 weeks and implying third-plus effects larger in absolute value than the physically feasible treatment dose. That pattern was the signature of a differential untreated path rather than a treatment effect, and it disappears once the control side is saturated.

\begin{table}[htbp]
\centering
\caption{Experience Gradient: Saturated Interaction Model (Ranking Points)}
\label{tab:saturated_experience}
\begin{threeparttable}
\scriptsize
\begin{tabular}{l*{5}{c}}
\toprule
 & 4w & 8w & 12w & 26w & 52w \\
\midrule
\multicolumn{6}{l}{\textit{Panel: GS-ATP}} \\ \midrule
$\hat{\beta}_h$ (first LL, $n=0$) & 55.2$^{***}$ & 58.4$^{***}$ & 69.1$^{***}$ & 61.4 & 4.7 \\
 & (18.5) & (18.7) & (20.1) & (41.1) & (64.4) \\
$\hat{\gamma}_{1h}$ (shift, $n=1$) & -6.7 & -5.6 & 3.7 & 55.4 & 70.4 \\
 & (38.2) & (34.9) & (37.4) & (60.6) & (108.5) \\
$\hat{\gamma}_{2h}$ (shift, $n\geq 2$) & -79.7$^{**}$ & -87.4$^{**}$ & -99.3$^{***}$ & -58.4 & -21.4 \\
 & (32.3) & (34.2) & (36.9) & (58.4) & (90.7) \\
Implied ATT, $n=1$ & 48.5 & 52.8 & 72.8 & 116.8 & 75.2 \\
Implied ATT, $n\geq 2$ & -24.5 & -29.0 & -30.2 & 2.9 & -16.7 \\
\\[0.3em]
\multicolumn{6}{l}{\textit{Panel: GS-WTA}} \\ \midrule
$\hat{\beta}_h$ (first LL, $n=0$) & 38.8 & 54.8$^{**}$ & 27.9 & 3.6 & 18.3 \\
 & (27.6) & (25.5) & (24.7) & (43.7) & (88.5) \\
$\hat{\gamma}_{1h}$ (shift, $n=1$) & 111.1$^{*}$ & 59.3 & -29.6 & -71.7 & -108.3 \\
 & (63.1) & (65.7) & (65.9) & (85.5) & (138.7) \\
$\hat{\gamma}_{2h}$ (shift, $n\geq 2$) & -61.4 & 6.0 & 20.0 & 128.7 & 82.7 \\
 & (96.9) & (93.5) & (102.3) & (114.9) & (186.5) \\
Implied ATT, $n=1$ & 149.9 & 114.1 & -1.7 & -68.0 & -90.0 \\
Implied ATT, $n\geq 2$ & -22.6 & 60.7 & 47.9 & 132.3 & 101.0 \\
\\[0.3em]
\multicolumn{6}{l}{\textit{Panel: NonGS-ATP}} \\ \midrule
$\hat{\beta}_h$ (first LL, $n=0$) & 4.2 & 6.0 & 1.0 & 34.7$^{**}$ & 28.9 \\
 & (8.4) & (9.6) & (10.1) & (14.3) & (28.6) \\
$\hat{\gamma}_{1h}$ (shift, $n=1$) & 13.3 & 14.4 & 18.6 & 5.4 & -14.1 \\
 & (11.0) & (13.7) & (15.2) & (23.8) & (46.9) \\
$\hat{\gamma}_{2h}$ (shift, $n\geq 2$) & 6.1 & 0.4 & 0.5 & -27.6 & -69.1$^{*}$ \\
 & (10.5) & (11.7) & (12.5) & (20.3) & (41.0) \\
Implied ATT, $n=1$ & 17.5 & 20.4 & 19.5 & 40.1 & 14.8 \\
Implied ATT, $n\geq 2$ & 10.3 & 6.4 & 1.4 & 7.1 & -40.3 \\
\\[0.3em]
\multicolumn{6}{l}{\textit{Panel: NonGS-WTA}} \\ \midrule
$\hat{\beta}_h$ (first LL, $n=0$) & 23.0$^{**}$ & 21.2$^{**}$ & 30.7$^{**}$ & 22.9 & 47.5 \\
 & (10.4) & (10.8) & (12.6) & (18.4) & (39.2) \\
$\hat{\gamma}_{1h}$ (shift, $n=1$) & -4.1 & -3.6 & -18.1 & -17.3 & -51.2 \\
 & (13.5) & (15.6) & (19.0) & (33.7) & (65.4) \\
$\hat{\gamma}_{2h}$ (shift, $n\geq 2$) & 17.4 & 4.8 & -13.0 & -67.8$^{**}$ & -152.6$^{**}$ \\
 & (13.7) & (16.8) & (19.0) & (30.6) & (68.2) \\
Implied ATT, $n=1$ & 19.0 & 17.5 & 12.6 & 5.6 & -3.8 \\
Implied ATT, $n\geq 2$ & 40.4 & 26.0 & 17.7 & -44.9 & -105.2 \\
\bottomrule
\end{tabular}

\begin{tablenotes}
\footnotesize
\item \textit{Notes.} Estimates of equation~(\ref{eq:saturated}) on the full unrestricted sample. $\hat\beta_h$ is the treatment effect for first-time recipients ($n = 0$); $\hat\gamma_{1h}$ and $\hat\gamma_{2h}$ shift that effect for second-time ($n = 1$) and third-plus ($n \geq 2$) recipients. Implied ATTs add the shifter to the base effect; their standard errors use the delta method and account for the covariance between $\hat\beta_h$ and $\hat\gamma_{jh}$. Grand Slam panels use Grand Slam $\times$ Year and horizon fixed effects; non-Grand-Slam panels add the control-function residual interacted with horizon. Pre-treatment controls include prior-LL counts, which vary within the pooled sample. Standard errors clustered at the player level. $^{***}p<0.01$, $^{**}p<0.05$, $^{*}p<0.10$.
\end{tablenotes}
\end{threeparttable}
\end{table}

\subsection{Verified Lottery Subsample and Contamination}
\label{sec:robustness_verified}

The Grand Slam lottery mechanism operates only when withdrawals occur before qualifying completes; post-qualifying withdrawals follow the ranking rule (the highest-ranked qualifying loser receives the slot). Withdrawal timing is not recorded in the data, so the main sample may pool lottery-assigned and ranking-assigned observations. This contamination has a signable direction: under ranking-based assignment, treated players are systematically higher-ranked than controls, biasing ranking-point effect estimates upward rather than toward attenuation. Two features of the analysis partially offset this bias. First, the pre-treatment control vector $Z^{pre}_{ie}$ includes ranking points (linear and quadratic), Elo, and surface Elo, absorbing the observable component of selection. Second, within-event fixed effects restrict comparisons to the top-four pool at a given tournament, where the ranking gap between treated and untreated candidates is mechanically small (at most three ranks).

To assess the residual contamination concern, I construct a \textit{verified lottery} subsample using two criteria, neither of which requires external information. Under the ranking rule an event that fills $k$ slots awards them to exactly the top $k$ ranked qualifying losers, in rank order, and each criterion identifies a distinct observable violation of that pattern.

The first is individual. A treated player whose rank among final-round qualifying losers exceeds the event's slot count cannot have been placed by the ranking rule, so their entry is provably a lottery draw. The second is at the event level. If $k$ slots were filled by rank order, the set of recipient ranks is exactly $\{1, 2, \ldots, k\}$; a gap in that set cannot arise under ranking and therefore proves that a lottery ran at that event. The second criterion identifies events the first misses: at a four-slot event whose recipients are ranked first and third, neither rank exceeds four, but the ranking rule cannot skip rank two. The gap test must be computed on the unrestricted sample rather than the first-time-restricted one, because repeat recipients excluded by the sample restriction would leave holes in the rank set that mimic genuine gaps. Both criteria share the assumption that no higher-ranked loser declined a slot, since a decline would also produce a gap.

Applying both criteria, 41 of the 64 Grand Slam events with first-time candidates are certified as lotteries, yielding 26 verified ATP events with $N = 64$ observations (35 treated, 29 control) and 15 verified WTA events with $N = 41$ observations (23 treated, 18 control).

Table~\ref{tab:verified_firstll} reports stacked dynamic estimates on this subsample. On the ATP tour the ranking-point coefficients keep the sign and the rising shape of the pooled estimates, running from 26.5 at 4 weeks to 87.5 at 26 weeks, but the standard errors are half again as large as in the pooled sample and only the 26-week estimate approaches significance ($p = 0.055$). On the WTA tour the 8-week gain survives at $+70.7$ ($p = 0.044$), close to the pooled estimate of 55.0. Elo is null at every horizon on both tours except a marginally negative ATP estimate at 52 weeks ($-51.3$, $p = 0.074$).

These results neither confirm nor refute the headline estimates, and the reason is power rather than contradiction. Table~\ref{tab:mde} reports minimum detectable effects: at 80\% power the verified ATP subsample detects roughly 90 to 125 ranking points over the 4-to-26-week range, against pooled point estimates of 72 to 83. The subsample is not capable of resolving an effect of the size the pooled sample identifies, so its nulls carry little information and elevating it to primary would substitute an underpowered null for a powered estimate. Nor can it corroborate them. The verified ATP point estimates at short horizons (26.5, 37.4, 41.6) run 50 to 64 percent below the pooled values, which is the direction upward contamination would produce; the differences are inside one verified standard error, so nothing is established either way, but the pattern should not be described as reassuring. A sample that cannot distinguish 72 from zero equally cannot distinguish 72 from 26.

An earlier and narrower version of this subsample, built on the individual rank criterion alone, produced a smaller sample ($N = 49$ ATP) on which Elo appeared to decline sharply at 12 and 52 weeks and WTA Elo appeared to rise at four of five horizons. Adding the event-level gap criterion enlarges the verified sample by roughly a third and those apparent effects attenuate to insignificance. The larger sample is the better-identified one, since the gap criterion certifies genuine lotteries that the narrower rule missed, so the attenuated estimates supersede the earlier ones. The episode illustrates how fragile inference is at these sample sizes and is the reason the paper rests its contamination defence on the placebo tests of Section~\ref{sec:robustness_placebo} rather than on this subsample.

The verified subsample is therefore uninformative in both directions, and the defence of the pooled estimates against contamination rests on the placebo tests of Section~\ref{sec:robustness_placebo} rather than on this subsample. Those tests target the contamination mechanism directly, since ranking-rule assignment selects on pre-treatment rank, and they do so at full sample size.

A partial news audit supports the classification rule, and its coverage should be stated precisely. Of the 65 first-time Grand Slam entries whose rank falls within the slot count, I searched tennis-media coverage and tournament records for 23: those ranked exactly first among qualifying losers, which is the subset where a ranking-rule assignment is most likely a priori, since that rule always fills the first vacancy with the top-ranked loser. Ten of the 23 were classified as ranking-rule assignments, seven of them with a specific withdrawal date and three on the basis of reported late-injury withdrawals without an exact date. Every one was a post-qualifying withdrawal, which triggers the ranking rule: Kyrgios withdrew from Wimbledon on the Sunday night before the 2023 main draw, Alexandrova from Wimbledon hours before her first-round match in 2024, Davidovich Fokina from Wimbledon the same year, Raonic from Roland Garros in 2021 and the US Open in 2019, Andreescu and Badosa from the US Open in 2023, and Tomljanovic from the Australian Open in 2023. None was a lottery draw. Thirteen were searched without yielding a datable source and remain unresolved; these are predominantly pre-2015 events with thin contemporaneous coverage. Appendix Table~\ref{tab:news_audit} lists every searched entry with the withdrawn player, the date where one could be established, and the classification reached.

The remaining 42 ambiguous entries, ranked second through fourth at events offering two to seven slots, have not been audited. They are classified as ambiguous on the rank criterion alone and are excluded from the verified subsample on that basis. This is the conservative treatment, because an unaudited entry is never counted as a verified lottery: auditing them can only move entries into the verified subsample or confirm their exclusion. The verified-lottery sample sizes reported above are therefore a lower bound on what complete verification would yield, and the corresponding estimates are correspondingly less precise than they would be under a full audit. Extending the audit to these 42 entries is the most direct available route to a larger verified subsample.

\begin{table}[htbp]
\centering
\caption{First-LL Verified Lottery Subsample: Dynamic Effects}
\label{tab:verified_firstll}
\begin{threeparttable}
\small
\begin{tabular}{l*{5}{c}}
\toprule
 & 4w & 8w & 12w & 26w & 52w \\
\midrule
\multicolumn{6}{l}{\textit{ATP}: $N = 49$ (24 treated, 25 control), 18 events} \\
  Ranking Pts $\Delta$ & 15.0 & 17.0 & 25.0 & 56.3 & -18.4 \\
  & (44.2) & (42.4) & (40.8) & (51.8) & (82.4) \\[0.3em]
  Elo $\Delta$ & -25.7$^{*}$ & -29.7$^{*}$ & -31.6$^{**}$ & 1.5 & -76.2$^{**}$ \\
  & (14.5) & (17.4) & (15.4) & (20.4) & (30.0) \\[0.3em]
  Main Draws & -0.36 & -0.63 & -0.34 & 0.66 & 1.15 \\
  & (0.89) & (0.81) & (0.81) & (0.93) & (2.26) \\[0.3em]
  Matches 250+ & -1.44 & -2.13 & -1.30 & 0.75 & 1.27 \\
  & (1.58) & (1.46) & (1.43) & (1.66) & (4.35) \\[0.3em]
\midrule
\multicolumn{6}{l}{\textit{WTA}: $N = 31$ (15 treated, 16 control), 10 events} \\
  Ranking Pts $\Delta$ & 64.2 & 81.5$^{**}$ & 34.0 & 46.3 & 129.2 \\
  & (53.5) & (39.4) & (40.2) & (67.8) & (146.7) \\[0.3em]
  Elo $\Delta$ & 21.9$^{**}$ & 35.6$^{***}$ & 32.2$^{**}$ & 23.5 & 59.7$^{**}$ \\
  & (8.9) & (11.0) & (13.4) & (16.6) & (26.7) \\[0.3em]
  Main Draws & 0.62 & 0.83 & 1.15 & 0.90 & 2.33 \\
  & (1.05) & (1.09) & (1.18) & (1.28) & (2.65) \\[0.3em]
  Matches 250+ & 0.48 & 1.81 & 1.92 & 2.79 & 4.27 \\
  & (1.91) & (1.74) & (1.78) & (1.99) & (4.91) \\[0.3em]
\midrule
\bottomrule
\end{tabular}

\begin{tablenotes}
\footnotesize
\item \textit{Notes.} Restricted to events with at least one LL entry whose rank among final-round qualifying losers exceeds the number of available LL slots at that event, which the ranking rule cannot produce. First-time LL candidates only. Stacked dynamic specification with Grand Slam $\times$ Year fixed effects, horizon fixed effects, and pre-treatment controls. Standard errors clustered at the player level. Minimum detectable effects for this subsample are reported in Table~\ref{tab:mde}. $^{***}p<0.01$, $^{**}p<0.05$, $^{*}p<0.10$.
\end{tablenotes}
\end{threeparttable}
\end{table}

\subsection{Pre-Treatment Balance}
\label{sec:robustness_balance}

Under genuine within-pool randomization, balance holds in expectation. Finite-sample imbalance is inevitable: with only four candidates per pool and a binary treatment, any single draw produces deviations that are properly interpreted as sampling variability rather than as evidence against randomization. Controls included in the outcome specification improve precision by reducing residual variance; they do not correct for a violation of random assignment, because under the lottery no such violation is present in expectation.

Table~\ref{tab:balance} reports formal balance tests on pre-treatment covariates, and the two tours look different enough that they warrant separate discussion.

On the ATP tour the picture is close to what randomization implies. The joint $F$-test yields $p = 0.106$, ranking points and Elo are balanced to three decimal places on the Elo measures, and the one covariate rejecting at the 5\% level is age ($p = 0.034$), where treated players average 24.9 years against 26.2 for controls, a gap of 1.3 years within a pool spanning 18 to 35. Age cannot plausibly be correlated with a lottery draw and enters $Z^{pre}_{ie}$ in any case. The within-event test yields $p = 0.046$, which with seven covariates across two tours is within the range chance produces.

The WTA tour is a different matter and the paper should not gloss it. Four of seven covariates reject at the 5\% level, and all four run in the same direction: treated players carry more ranking points (0.53 against 0.42, $p = 0.007$; the squared term $p = 0.006$) and a higher Elo (18.54 against 17.42, $p = 0.022$; squared term $p = 0.022$). The two pairs are collinear by construction, so this is roughly two independent dimensions rather than four, but both are precisely the dimensions on which ranking-rule assignment would select. The joint $F$-test does not reject at conventional levels ($p = 0.063$), and the within-event test is far from rejecting ($p = 0.516$), which suggests the imbalance operates between events rather than within the pools where the lottery runs. That is the reading most consistent with the design, since event fixed effects absorb between-pool variation, and it is supported by the pre-treatment placebo tests of Section~\ref{sec:robustness_placebo}, which are null on the WTA tour at all three horizons. It is not, however, a reading the balance table alone establishes. The honest position is that the WTA lottery sample shows level imbalance on exactly the covariates contamination would move, that the within-event and placebo evidence argues against this reflecting non-random assignment within pools, and that the WTA estimates should be read with more caution than the ATP estimates on that account. Section~\ref{sec:results_dynamics} already treats the WTA ranking-point result as suggestive rather than established on independent grounds, and this balance evidence reinforces that reading.

Event study plots (Figures~\ref{fig:event_study_atp} and~\ref{fig:event_study_wta}) show parallel trajectories in the weeks preceding each Grand Slam. Non-GS balance tests are reported in the appendix; that sample is not randomized, so imbalance there is expected and motivates the control function.

\begin{table}[htbp]
\centering
\caption{Pre-Treatment Balance: Grand Slam Samples}
\label{tab:balance}
\begin{threeparttable}
\small
\begin{tabular}{l rrr rrr}
\toprule
 & \multicolumn{3}{c}{ATP} & \multicolumn{3}{c}{WTA} \\
\cmidrule(lr){2-4} \cmidrule(lr){5-7}
Variable & LL & Control & $p$ & LL & Control & $p$ \\
\midrule
Ranking points / 1000 & 0.48 & 0.45 & 0.088 & 0.53 & 0.42 & 0.007 \\
(Ranking points / 1000)$^2$ & 0.25 & 0.21 & 0.083 & 0.31 & 0.20 & 0.006 \\
Elo / 100 & 17.30 & 17.30 & 0.977 & 18.54 & 17.42 & 0.022 \\
(Elo / 100)$^2$ & 299.27 & 299.28 & 0.977 & 343.91 & 303.37 & 0.022 \\
Surface Elo / 100 & 16.29 & 16.38 & 0.686 & 17.62 & 17.28 & 0.407 \\
(Surface Elo / 100)$^2$ & 266.91 & 270.07 & 0.676 & 314.28 & 301.26 & 0.366 \\
Age & 24.93 & 26.20 & 0.034 & 23.15 & 23.83 & 0.391 \\
\midrule
Joint $F$-test $p$-value & \multicolumn{3}{c}{0.106} & \multicolumn{3}{c}{0.063} \\
$N$ & \multicolumn{3}{c}{120} & \multicolumn{3}{c}{82} \\
\bottomrule
\end{tabular}

\begin{tablenotes}
\footnotesize
\item \textit{Notes.} Columns compare means for treated (LL selected) and control (unselected) players at the player-episode level. Covariates include ranking points / 1000 (linear and quadratic), Elo / 100 (linear and quadratic), surface Elo / 100 (linear and quadratic), and age. Prior LL variables excluded (zero by construction for first-time candidates). Joint $F$-statistic tests the null that all covariates are jointly balanced. First-time ATP: $N = 120$; first-time WTA: $N = 82$. Grand Slam sample, 2006--2024.
\end{tablenotes}
\end{threeparttable}
\end{table}

\subsection{Pre-Treatment Placebo Tests}
\label{sec:robustness_placebo}

Balance on levels and visually parallel event-study trajectories are suggestive, but the direct test of random assignment is whether eventual treatment status predicts outcomes measured \textit{before} treatment. I construct ranking-point changes over the twelve, eight, and four weeks preceding each focal event, using the same range-join to the weekly ranking files that produces the post-treatment outcomes, and regress each on eventual LL status with the headline specification.

Table~\ref{tab:placebo} reports the results. None of the twelve tests, spanning three pre-treatment horizons, two tours, and both the pooled and verified-lottery samples, rejects at the 5\% level; the largest absolute $t$-statistic in the set is 1.47. Point estimates are small relative to the post-treatment effects: on the pooled ATP sample they run from $-4.3$ to $+5.9$ ranking points, against post-treatment effects of 72 to 83, and they alternate in sign rather than trending. This is the sharpest available evidence against the residual contamination concern. Ranking-rule assignment selects the highest-ranked qualifying loser, so contaminated observations would carry pre-treatment ranking trajectories that differ systematically between treated and control players. No such difference appears.

\begin{table}[htbp]
\centering
\caption{Pre-Treatment Placebo Tests}
\label{tab:placebo}
\begin{threeparttable}
\small
\begin{tabular}{lll*{3}{c}}
\toprule
Sample & Tour & & $-12$w & $-8$w & $-4$w \\
\midrule
Pooled & ATP & $\hat{\beta}$ & -3.0 & 5.9 & -4.3 \\
 & & & (20.7) & (15.5) & (10.3) \\[0.3em]
Pooled & WTA & $\hat{\beta}$ & 40.0 & 22.1 & -15.0 \\
 & & & (27.1) & (20.8) & (11.8) \\[0.3em]
Verified & ATP & $\hat{\beta}$ & -9.2 & -3.6 & -4.6 \\
 & & & (32.6) & (22.8) & (19.0) \\[0.3em]
Verified & WTA & $\hat{\beta}$ & -36.3 & -13.4 & -11.6 \\
 & & & (56.1) & (45.0) & (24.3) \\[0.3em]
\bottomrule
\end{tabular}

\begin{tablenotes}
\footnotesize
\item \textit{Notes.} Outcome is the change in ranking points from the indicated pre-treatment week to the focal event date, so all variation predates assignment. Specification matches the headline dynamic model with Grand Slam $\times$ Year fixed effects and pre-treatment controls. Standard errors clustered at the player level. Under random assignment these coefficients should be indistinguishable from zero. $^{***}p<0.01$, $^{**}p<0.05$, $^{*}p<0.10$.
\end{tablenotes}
\end{threeparttable}
\end{table}

\subsection{Attrition and Lee Bounds}
\label{sec:robustness_attrition}

Players leave the rankings, and outcomes are observed only for those who remain. If treatment itself keeps players ranked, long-horizon outcomes condition on a treatment-affected selection event. Table~\ref{tab:attrition} reports observation rates. Through 26 weeks the rates are close: 80.3\% of treated and 78.0\% of control ATP players are observed at 26 weeks, a gap of 2.4 percentage points that a two-proportion test cannot distinguish from zero ($p = 0.93$). At 52 weeks the gap widens to 12.7 percentage points on the ATP tour (75.4\% versus 62.7\%), and reverses to $-10.0$ points on the WTA tour; neither is individually significant at conventional levels given the sample sizes ($p = 0.19$ and $p = 0.39$), but both are large enough to matter for interpretation.

Table~\ref{tab:lee} reports \citet{Lee2009_bounds} bounds, which trim the over-observed group from each tail of its outcome distribution to equalize observation rates and so bound the effect for the subpopulation observed under either assignment. Through 26 weeks the bounds are tight and exclude zero on the ATP tour: at 12 weeks the unadjusted mean difference of 44.4 ranking points is bounded within $[38.3, 50.1]$, and at 26 weeks within $[23.4, 46.8]$. The short- and medium-horizon results are therefore not an artifact of differential survival. At 52 weeks the bounds widen sharply and span zero on both tours ($[-143.9, 35.3]$ ATP; $[-18.2, 90.7]$ WTA), because the trimming fraction rises to 17\% and 12\% respectively. The paper's fade-out claim at 52 weeks is consequently the one result that differential attrition can also explain, and Section~\ref{sec:conclusion} reports it with that caveat rather than as established.

\begin{table}[htbp]
\centering
\caption{Attrition: Outcome Observation Rates by Treatment Status}
\label{tab:attrition}
\begin{threeparttable}
\small
\begin{tabular}{ll*{5}{c}}
\toprule
Tour & & 4w & 8w & 12w & 26w & 52w \\
\midrule
ATP & Treated (\%) & 100.0 & 91.8 & 86.9 & 80.3 & 75.4 \\
 & Control (\%) & 100.0 & 94.9 & 89.8 & 78.0 & 62.7 \\
 & Difference & 0.0 & -3.1 & -2.9 & 2.4 & 12.7 \\[0.3em]
WTA & Treated (\%) & 100.0 & 97.0 & 97.0 & 93.9 & 75.8 \\
 & Control (\%) & 100.0 & 100.0 & 98.0 & 91.8 & 85.7 \\
 & Difference & 0.0 & -3.0 & -1.0 & 2.1 & -10.0 \\[0.3em]
\bottomrule
\end{tabular}

\begin{tablenotes}
\footnotesize
\item \textit{Notes.} Share of players with a non-missing ranking-point outcome at each horizon. Difference row reports treated minus control in percentage points, with significance from a two-proportion test. First-time Grand Slam samples. $^{***}p<0.01$, $^{**}p<0.05$, $^{*}p<0.10$.
\end{tablenotes}
\end{threeparttable}
\end{table}

\begin{table}[htbp]
\centering
\caption{Lee (2009) Bounds on the Ranking-Point Effect}
\label{tab:lee}
\begin{threeparttable}
\small
\begin{tabular}{ll*{5}{c}}
\toprule
Tour & & 4w & 8w & 12w & 26w & 52w \\
\midrule
ATP & Unadjusted & 22.4 & 27.7 & 44.4 & 37.2 & -42.0 \\
 & Lee lower & 22.4 & 23.4 & 38.3 & 23.4 & -143.9 \\
 & Lee upper & 22.4 & 32.3 & 50.1 & 46.8 & 35.3 \\[0.3em]
WTA & Unadjusted & 23.8 & 35.6 & 10.8 & -11.5 & 19.2 \\
 & Lee lower & 23.8 & 32.4 & 10.8 & -11.5 & -18.2 \\
 & Lee upper & 23.8 & 44.1 & 10.8 & -11.5 & 90.7 \\[0.3em]
\bottomrule
\end{tabular}

\begin{tablenotes}
\footnotesize
\item \textit{Notes.} Unadjusted row reports the raw difference in mean ranking-point change between treated and control players among those observed. Lee bounds trim the over-observed group from each tail of its outcome distribution until observation rates are equal, bounding the effect for the always-observed subpopulation under monotone selection. Bounds coincide with the unadjusted difference where observation rates are already equal. First-time Grand Slam samples.
\end{tablenotes}
\end{threeparttable}
\end{table}

\subsection{Multiple Testing}
\label{sec:robustness_mht}

The paper reports four outcomes at five horizons on four samples, and the resulting 80 tests invite a correction. Which correction is appropriate depends on what is being controlled, and the three standards reported here give materially different answers, so I state all of them rather than selecting the most favourable.

Within the ranking-point outcome, treating the five horizons as one family of 20 tests, the ATP Grand Slam effects at 4, 8, and 12 weeks survive Benjamini-Hochberg correction ($q = 0.039$, $0.039$, $0.020$); the 26-week effect does not ($q = 0.119$), and the WTA 8-week effect is marginal ($q = 0.067$). Applying Benjamini-Hochberg jointly across all 80 tests gives a similar picture: 23 tests survive at $q < 0.05$, including the three ATP Grand Slam horizons ($q = 0.028$, $0.028$, $0.009$).

Romano-Wolf step-down inference is stricter and gives a different answer. Where Benjamini-Hochberg bounds the expected share of rejections that are false, Romano-Wolf bounds the probability of making any false rejection at all, and it does so using a bootstrap that preserves the dependence among the test statistics. Of the 80 tests, 14 survive at the 5\% level, and all 14 are non-Grand-Slam access outcomes: main draws entered and matches at 250-level events on both tours, where the raw $p$-values fall below $0.001$. No Grand Slam ranking-point effect survives. On the ATP tour the step-down $p$-values are $0.685$ at 4 weeks, $0.685$ at 8 weeks, and $0.410$ at 12 weeks; the WTA 8-week effect returns $0.815$.

The right way to read this is not that the Grand Slam results are spurious, but that family-wise error control across 80 tests is a demanding standard that a single design with 120 observations was never going to meet. The dynamic path across horizons is one economic claim rather than five independent ones, and the appropriate correction for a pre-specified primary outcome is the within-outcome family, which the Grand Slam effects clear. What the Romano-Wolf results do establish is a ranking of evidential strength. The non-Grand-Slam access effects are robust to the strictest available correction and are the most secure findings in the paper. The Grand Slam ranking-point effects are supported by randomization inference, by placebo tests, and by attrition bounds, and they survive false-discovery-rate control, but they are not individually strong enough to clear family-wise control over the full outcome set. Readers who require the latter standard should treat the Grand Slam estimates as suggestive and the non-Grand-Slam access results as established.

\subsection{Control Function Specification Check}
\label{sec:robustness_cf}

The CF specification check (Appendix Table~\ref{tab:cf_validation}) applies the control function to the randomized GS sample using a uniform selection probability $P_i^{LL} = c_e / N_e$. On the WTA tour, the specification passes: $\hat{\rho}_h$ is insignificant at all horizons for ranking points (e.g., $\hat{\rho}_4 = 17.5$, SE $48.5$) and treatment effect estimates change modestly. On the ATP tour, the specification fails: $\hat{\rho}_h$ is large and highly significant (e.g., $\hat{\rho}_4 = 90.6$, $p < 0.01$), indicating that the uniform probability is a poor approximation for the ATP GS assignment process. This likely reflects that the ATP GS pool size varies substantially across events (from 4 to 10+ candidates depending on the qualifying draw size and withdrawal timing), so the uniform-within-top-4 assumption mischaracterizes the selection. The ATP GS failure does not compromise the lottery-based estimates, which rely only on random assignment within the top-four pool rather than on a functional form assumption. It does indicate that the non-GS CF estimates should be read as corroboration rather than as a free-standing identification strategy, and that the lottery design is the primary basis for causal inference.

\subsection{Wildcards}
\label{sec:robustness_wildcards}

Wildcard entries provide a parallel mechanism for main-draw access. If control players differentially receive wildcards to subsequent events, the counterfactual is contaminated. In the first-LL sample, wildcards are rare among qualifying-level players (ranked 100th--300th), who are generally below the threshold for wildcard consideration at tour-level events. In the 26-week window following the focal event, the average number of wildcard main-draw entries is 0.08 per treated player and 0.09 per control player on the ATP tour ($p = 0.71$), and 0.12 versus 0.11 on the WTA tour ($p = 0.83$). Differential wildcard receipt is therefore too small to confound the LL treatment effect.

\subsection{Small-Sample Inference}
\label{sec:robustness_power}

The Grand Slam samples are small, so asymptotic cluster-robust inference warrants scrutiny. Table~\ref{tab:clusters} reports observation counts, treated and control counts, player clusters, and events for every sample. Because players rarely appear in the pool more than once under the first-time restriction, cluster counts track sample sizes closely: 109 players in the pooled ATP sample, 77 in the pooled WTA sample, and 62 and 40 in the verified-lottery ATP and WTA subsamples. The verified WTA subsample sits at the boundary of the range where asymptotic clustered inference is normally considered reliable, so results on both verified subsamples are reported alongside bootstrap $p$-values below.

Table~\ref{tab:mde} reports minimum detectable effects at 80\% power and a 5\% two-sided test, computed as $(z_{0.975} + z_{0.80})$ times the realised clustered standard error. These convert null results into statements about what the data can and cannot rule out, and the comparison across samples is stark. On the pooled ATP sample the minimum detectable ranking-point effect is 74.8 at 4 weeks, 74.4 at 8 weeks, and 73.1 at 12 weeks, bracketing the estimated effects of 72.5, 72.1, and 82.8. At 4 and 8 weeks the minimum detectable effect slightly exceeds the estimate; only at 12 weeks does it fall below. The design has approximately the power required to detect effects of the size it reports, and no margin beyond that. On the verified-lottery ATP subsample the minimum detectable effect rises to 92 to 125 over the same range while the point estimates fall, so a null there is close to uninformative. The pattern is sharper still on the WTA side, where the verified subsample requires 95 to 129 points at short horizons.

Two implications follow. First, the paper is operating at the edge of its power throughout, and results should be read with the corresponding caution: this is a design with 120 and 82 observations, not one that can resolve modest effects. Second, and for the same reason, the null results on the verified subsample carry little evidential weight against the pooled estimates, since that subsample could not have detected the pooled effect had it been present.

Table~\ref{tab:wildboot} reports wild cluster bootstrap $p$-values following \citet{Cameron2008_wildbootstrap}, using the restricted bootstrap with the null imposed and Webb six-point weights, which outperform Rademacher weights when the cluster count is small. The bootstrap re-estimates the full specification 1{,}999 times per coefficient.

The bootstrap corroborates the asymptotic inference rather than overturning it, and on the pooled samples it is if anything more favourable. On the pooled ATP sample the bootstrap $p$-values at 4, 8, and 12 weeks fall below 0.001 against asymptotic values of 0.008, 0.008, and 0.002, and the 26-week bootstrap $p$-value is 0.026 against an asymptotic 0.047. On the pooled WTA sample the 8-week bootstrap $p$-value is 0.012 against 0.017. No coefficient that is significant under clustered inference loses significance under the bootstrap, on any sample or at any horizon. Cluster-robust standard errors are therefore not over-rejecting here, which is consistent with the cluster counts in Table~\ref{tab:clusters} being larger than the cases where the small-cluster problem is acute.

The 26-week ATP effect is the one place where the inference methods genuinely disagree, and the disagreement is informative about which to trust. It is significant under clustered standard errors ($p = 0.047$) and under the wild bootstrap ($p = 0.026$), but not under Fisher randomization inference ($p_{\text{Fisher}} = 0.103$). The two model-based procedures share an assumption that randomization inference does not: they treat the regression as correctly specified and resample from it, whereas the Fisher test permutes the treatment indicator within the blocks in which it was actually assigned and so tests the sharp null using the design itself. For a lottery, randomization inference is the more appropriate standard, and the paper accordingly treats the 26-week result as the weakest of the significant estimates despite two procedures endorsing it.

\begin{table}[htbp]
\centering
\caption{Sample and Cluster Counts}
\label{tab:clusters}
\begin{threeparttable}
\small
\begin{tabular}{lrrrrr}
\toprule
Sample & $N$ & Treated & Control & Players & Events \\
\midrule
GS-ATP (pooled first-LL) & 120 & 61 & 59 & 109 & 55 \\
GS-WTA (pooled first-LL) & 82 & 33 & 49 & 77 & 32 \\
GS-ATP (verified lottery) & 64 & 35 & 29 & 62 & 26 \\
GS-WTA (verified lottery) & 41 & 23 & 18 & 40 & 15 \\
NonGS-ATP (first-LL) & 2054 & 523 & 1531 & 830 & 513 \\
NonGS-WTA (first-LL) & 1423 & 378 & 1045 & 736 & 422 \\
\bottomrule
\end{tabular}

\begin{tablenotes}
\footnotesize
\item \textit{Notes.} Player clusters are the unit at which standard errors are clustered throughout. Events count distinct tour $\times$ tournament $\times$ year combinations contributing identifying variation.
\end{tablenotes}
\end{threeparttable}
\end{table}

\begin{table}[htbp]
\centering
\caption{Minimum Detectable Effects (Ranking Points, 80\% Power)}
\label{tab:mde}
\begin{threeparttable}
\small
\begin{tabular}{ll*{5}{c}}
\toprule
Sample & & 4w & 8w & 12w & 26w & 52w \\
\midrule
Pooled-ATP & $\hat{\beta}_h$ & 72.5 & 72.1 & 82.8 & 72.6 & 23.5 \\
 & MDE (80\%) & 74.8 & 74.4 & 73.1 & 101.4 & 165.5 \\[0.3em]
Verified-ATP & $\hat{\beta}_h$ & 15.0 & 17.0 & 25.0 & 56.3 & -18.4 \\
 & MDE (80\%) & 123.9 & 118.8 & 114.4 & 145.2 & 230.8 \\[0.3em]
Pooled-WTA & $\hat{\beta}_h$ & 40.0 & 55.0 & 27.8 & 7.5 & 26.5 \\
 & MDE (80\%) & 79.9 & 63.0 & 56.4 & 119.9 & 242.0 \\[0.3em]
Verified-WTA & $\hat{\beta}_h$ & 64.2 & 81.5 & 34.0 & 46.3 & 129.2 \\
 & MDE (80\%) & 150.0 & 110.4 & 112.5 & 190.1 & 411.0 \\[0.3em]
\bottomrule
\end{tabular}

\begin{tablenotes}
\footnotesize
\item \textit{Notes.} Minimum detectable effect at 80\% power for a 5\% two-sided test, computed as $2.80 \times \text{SE}$ using the realised player-clustered standard error from the stacked dynamic specification. A null result is informative only where the minimum detectable effect lies below the effect size of interest.
\end{tablenotes}
\end{threeparttable}
\end{table}

\begin{table}[htbp]
\centering
\caption{Wild Cluster Bootstrap Inference (Ranking Points)}
\label{tab:wildboot}
\begin{threeparttable}
\small
\begin{tabular}{ll*{5}{c}}
\toprule
Sample & & 4w & 8w & 12w & 26w & 52w \\
\midrule
Verified-ATP & $\hat{\beta}_h$ & 26.5 & 37.4 & 41.6 & 87.5 & 29.1 \\
 & $p$ (clustered) & 0.444 & 0.279 & 0.211 & 0.055 & 0.675 \\
 & $p$ (wild boot.) & 0.481 & 0.334 & 0.286 & 0.072 & 0.681 \\[0.3em]
Verified-WTA & $\hat{\beta}_h$ & 67.2 & 70.7 & 24.1 & -19.5 & -24.2 \\
 & $p$ (clustered) & 0.151 & 0.044 & 0.498 & 0.771 & 0.850 \\
 & $p$ (wild boot.) & 0.162 & 0.040 & 0.464 & 0.804 & 0.858 \\[0.3em]
Pooled-ATP & $\hat{\beta}_h$ & 72.5 & 72.1 & 82.8 & 72.6 & 23.5 \\
 & $p$ (clustered) & 0.008 & 0.008 & 0.002 & 0.047 & 0.692 \\
 & $p$ (wild boot.) & 0.000 & 0.000 & 0.000 & 0.026 & 0.665 \\[0.3em]
Pooled-WTA & $\hat{\beta}_h$ & 40.0 & 55.0 & 27.8 & 7.5 & 26.5 \\
 & $p$ (clustered) & 0.165 & 0.017 & 0.171 & 0.861 & 0.760 \\
 & $p$ (wild boot.) & 0.151 & 0.012 & 0.171 & 0.887 & 0.757 \\[0.3em]
\bottomrule
\end{tabular}

\begin{tablenotes}
\footnotesize
\item \textit{Notes.} Restricted wild cluster bootstrap with the null imposed, Webb six-point weights, and 1{,}999 replications, clustered at the player level \citep{Cameron2008_wildbootstrap}. Clustered $p$-values are the asymptotic values reported elsewhere in the paper; bootstrap $p$-values are the share of bootstrap $t$-statistics exceeding the observed value in absolute terms.
\end{tablenotes}
\end{threeparttable}
\end{table}
